% Part: proof-theory
% Chapter: normalization
% Section: introduction
%READERNOTE{SOL6-A307: उपसूत्राच्या गुणधर्माची परिचयातील सूचना पुढील अचूक निष्कर्षाशी जुळवली: उघडी गृहीतके, संख्यापकांची पद-उदाहरणे आणि आवश्यक असत्य समाविष्ट आहेत.}

\documentclass[../../../include/open-logic-section]{subfiles}

\begin{document}

\olfileid{pt}{nor}{int}

\olsection{परिचय}

~$!A \lif !B$ ही शर्त तुम्ही
!!{prove}d केली असेल, तर
\Elim{\lif} वापरून~$!B$ ची
!!a{proof} करताना तिचा उपयोग
होऊ शकतो. अर्थात त्यासाठी~$!A$
ची !!a{proof}~$\delta_1$ ही हवीच.
~$!A \lif !B$ ची तुमची
!!{proof}~$\delta_3$ बहुधा
\Intro{\lif} ने संपते;
तो नियम उघड्या गृहीतक~$!A$
पासून~$!B$ ची
!!a{proof}~$\delta_2$
घेऊन~$!A \lif !B$
ची सिद्धता बनवतो.
तसे असल्यास तुमची अंतिम
!!{proof} पुढील रूपाची असते:
\[
\AxiomC{$\Discharge{!A}{x}$}
\RightLabel{$\delta_2$}
\DeduceC{$!B$}
\DischargeRule{\Intro{\lif}}{x}
\UnaryInfC{$!A \lif !B$}
\AxiomC{}
\RightLabel{$\delta_1$}
\DeduceC{$!A$}
\RightLabel{\Elim{\lif}}
\BinaryInfC{$!B$}
\DisplayProof
\]
आधीच मिळालेली
$!A \lif !B$ ची !!a{proof}
पुन्हा वापरली म्हणून
तुमची पद्धत कार्यक्षम असली,
तरी !!{proof} स्वतः
कार्यक्षम नाही.
$!A \lif !B$ आणून लगेच
निरसित करण्याचा हा
वळसा वाटतो. ~$!B$
!!{prove} करण्याचा अधिक
थेट मार्ग असायला हवा.

तसा मार्ग नेहमीच असतो.
अशा “वळशांना”—!!a{introduction}
नियमाच्या लगेच नंतर
!!a{elimination} नियम लावल्यास
तयार होणाऱ्या रचनांना—
नैसर्गिक निगमनाच्या
!!{proof}s मधून नेहमी
काढून टाकता येते.
याला \emph{सामान्यरूपीकरण}
म्हणतात; परिणामी मिळणारी
!!{proof} ही \emph{सामान्यरूप}
असते (किंवा ती !!a{proof}
\emph{सामान्य रूपात} आहे
असे म्हणतात).

या निष्कर्षाची सिद्धता
क्रमवर्ती कलनातील
कट-निर्मूलन प्रमेयाप्रमाणेच
काहीशी गुंतागुंतीची आहे;
तिच्यात अनेक प्रसंग
पाहावे लागतात.
अंतःप्रज्ञावादी~\Log{N1i}
पेक्षा अभिजात~\Log{N1c}
साठी ती सिद्ध करणे
अधिक कठीण असते.
क्रमवर्ती कलनात~\Cut{}
नियम वापरल्यामुळे विना-कट
!!{prove} करता येत नाही असे
अधिक काही सिद्ध करता येत
नाही, हे कट-निर्मूलन दाखवते.
त्याचप्रमाणे वळसे टाळून
मिळू शकत नाहीत असे
अधिक निष्कर्ष वळसे वापरूनही
मिळत नाहीत, हे
सामान्यरूपीकरण दाखवते.
इतर साम्येही आहेत:
एखाद्या निष्कर्षाची
वळसांसह सर्वात लहान
!!{proof} जितकी असेल,
त्याच निष्कर्षाची
सामान्यरूप !!{proof}
त्याहून खूप मोठी असू शकते.
तसेच क्रमवर्ती कलनात
विना-कट सर्वात लहान
!!{proof} पेक्षा
कट असलेली !!{proof}s
खूप लहान असू शकतात.
उप-!!{formula} गुणधर्मही
यातून मिळतो: नैसर्गिक
निगमनाच्या !!{proof} मध्ये
निष्कर्ष आणि उघडी गृहीतके
यांच्या उप-!!{formula}s, संख्यापकांमुळे लागणारी
त्यांची पद-आदेशनाची उदाहरणे आणि आवश्यक असत्य
वापरून निष्कर्ष नेहमी
!!{prove} करता येतो.
क्रमवर्ती कलनात हा गुणधर्म
कट-निर्मूलनातून मिळतो,
तसाच येथे तो
सामान्यरूपीकरणातून मिळतो.

\end{document}
